\documentclass[11pt]{article}
\usepackage[hyperref]{acl2023}
\usepackage{times}
\usepackage{latexsym}
\usepackage{amsmath}
\usepackage{booktabs}
\usepackage{graphicx}
\usepackage{multirow}
\usepackage{xcolor}
\usepackage{enumitem}

\title{Gauntlet: Defense-in-Depth Prompt Injection Detection\\via NLI-Augmented Cascade}

\author{Ashwin Shanmugasundaram \\
  Department of Computer Engineering \\
  Virginia Tech \\
  \texttt{ashwin@vt.edu}}

\begin{document}
\maketitle

\begin{abstract}
Prompt injection is a critical vulnerability in LLM-powered applications, where adversarial inputs manipulate a model into ignoring its system instructions. Existing open-source detectors typically rely on a single classification model, leading to either high false positive rates or low recall. We present Gauntlet, an open-source, fully local prompt injection detector that uses a three-layer defense-in-depth cascade: rule-based pattern matching, embedding similarity, and a novel DeBERTa-v3-base ensemble. Our key contribution is reframing prompt injection detection as a Natural Language Inference (NLI) task---asking whether input text \textit{entails} the hypothesis ``This text attempts to override system instructions''---rather than standard binary classification. A pre-trained NLI model achieves 9.7\% false positive rate on the NotInject over-defense benchmark with zero task-specific training, compared to 20.9\% for our trained binary classifier, demonstrating that task framing fundamentally alters the decision boundary. By combining the binary and NLI classifiers in a both-agree ensemble, Gauntlet achieves 4.1\% over-defense FPR (95\% CI: [2.1\%, 6.5\%]) and 0.701 F1 on ProtectAI-Validation (95\% CI: [0.679, 0.721]), significantly outperforming Meta Prompt Guard 2 ($p < 0.0001$, McNemar's test) on general detection while running entirely on consumer hardware (RTX 3060, 6GB VRAM). Code, models, and evaluation scripts are publicly available under Apache 2.0.
\end{abstract}

% ============================================================================
\section{Introduction}

Large Language Model (LLM) applications are increasingly deployed in production settings---customer service agents, code assistants, document processors---where they receive untrusted user input alongside trusted system prompts. Prompt injection attacks exploit this trust boundary: an adversary crafts input that causes the model to disregard its system instructions and execute attacker-controlled directives instead \cite{schulhoff2023hackaprompt, greshake2023indirect}.

Several open-source prompt injection detectors exist, including Meta Prompt Guard 2 \cite{meta2025pg2}, ProtectAI's DeBERTa v2, and the deepset classifier. However, each uses a single classification model, creating a fundamental tension: high recall requires aggressive detection thresholds, which increases false positives on legitimate inputs. This is particularly problematic for benign text containing security-adjacent vocabulary---phrases like ``ignore the formatting issues'' or ``override the default settings''---which single-model detectors frequently misclassify \cite{lee2025injecguard}.

We address this with three contributions:

\begin{enumerate}[leftmargin=*]
\item \textbf{NLI framing for prompt injection detection.} We show that reframing detection as a Natural Language Inference task---determining whether user input entails the hypothesis ``This text attempts to override system instructions''---reduces false positives by 55\% compared to binary classification, even in a zero-shot setting with no task-specific training. This finding suggests that the task formulation matters more than model scale for over-defense mitigation.

\item \textbf{A both-agree ensemble strategy.} We combine a binary classifier (high precision) with an NLI classifier (high recall) and require both models to agree before flagging an input. This eliminates the majority of false positives from either model alone, achieving 4.1\% FPR on the NotInject over-defense benchmark---competitive with or better than all evaluated open-source detectors.

\item \textbf{A defense-in-depth cascade architecture.} Our three-layer cascade (regex rules, embedding similarity, DeBERTa ensemble) routes easy cases to cheap layers, reserving expensive model inference for ambiguous inputs. This follows OWASP and NIST defense-in-depth recommendations and enables sub-millisecond detection for 47\% of attacks on ProtectAI-Validation.
\end{enumerate}

Gauntlet runs entirely locally on consumer hardware (NVIDIA RTX 3060, 6GB VRAM), requires no cloud API calls, and is available on PyPI (\texttt{pip install gauntlet-ai}) under Apache 2.0.

% ============================================================================
\section{Related Work}

\paragraph{Single-model detectors.} Meta Prompt Guard 2 \cite{meta2025pg2} uses mDeBERTa-base with energy-based out-of-distribution loss \cite{liu2020energy} and an adversarial tokenizer, reporting 97.5\% recall at 1\% FPR on a private benchmark across 8 languages. ProtectAI's DeBERTa v2, the most downloaded open-source detector (419K+ monthly downloads), achieves high recall (99.7\%) but suffers from high false positive rates. The deepset classifier, a pioneer in open-source prompt injection detection, was trained on only 662 samples and shows limited generalization. PIGuard \cite{lee2025injecguard} uses DeBERTa-v3-base with MOF (Mitigating Over-defense for Free) training, achieving the best reported over-defense mitigation among single-model approaches.

\paragraph{Over-defense as a first-class metric.} Lee et al.~\cite{lee2025injecguard} introduced NotInject, a benchmark of 339 benign samples containing trigger words (e.g., ``ignore,'' ``override,'' ``system prompt'') at three difficulty levels. InjecGuard formalized three-dimensional evaluation: malicious accuracy, benign accuracy, and over-defense accuracy. Our work follows this paradigm, treating false positive rate on trigger-word-containing benign text as a primary evaluation metric.

\paragraph{NLI for text classification.} Yin et al.~\cite{yin2019benchmarking} demonstrated that NLI models can perform zero-shot text classification by testing entailment between input text and task-descriptive hypotheses. Laurer et al.~\cite{laurer2023building} showed that NLI-based zero-shot classification generalizes across domains with minimal performance loss. To our knowledge, we are the first to apply NLI framing specifically to prompt injection detection.

\paragraph{Evaluation concerns.} Polyakov et al.~\cite{polyakov2026benchmarks} demonstrated that standard train/test evaluation of prompt injection classifiers inflates performance by approximately 8.4 AUC points compared to Leave-One-Dataset-Out (LODO) evaluation, due to dataset-specific shortcuts. We acknowledge this limitation and report results on four independent benchmarks, including two (deepset, JailbreakBench) whose data was not used in training.

% ============================================================================
\section{Method}

\subsection{Architecture Overview}

Input passes through layers sequentially; detection at any layer halts the cascade. This design ensures that obvious attacks (caught by regex) never incur model inference costs, while subtle attacks are evaluated by the full ensemble. The cascade comprises: (1) adversarial preprocessing, (2) regex-based pattern matching, (3) embedding similarity, and (4) a DeBERTa binary+NLI ensemble.

\subsection{Adversarial Preprocessing}

Before any detection layer, all input undergoes adversarial text sanitization. This step strips characters that are invisible to humans but can manipulate tokenizer behavior: zero-width characters (U+200B, U+200C/D, U+FEFF, U+2060--U+2064), Unicode Tags block (U+E0001--U+E007F, used for ASCII smuggling attacks achieving 90--100\% attack success rates), variation selectors (U+FE00--U+FE0F), and directional overrides (U+202A--U+202E). The preprocessor normalizes exotic whitespace to ASCII space and collapses whitespace runs, preserving legitimate CJK, Arabic, Hindi, emoji, and accented Latin text. We note that DeBERTa-v3's SentencePiece (Unigram) tokenizer is inherently resistant to TokenBreak attacks that affect BPE-based models.

\subsection{Layer 1: Rule-Based Detection}

Layer 1 applies 50+ regex patterns organized across 9 attack categories: instruction override, jailbreak, delimiter injection, data extraction, context manipulation, obfuscation, hypothetical framing, multilingual injection, and indirect injection. Patterns cover 13 languages and include NFKC normalization and a homoglyph table mapping $\sim$80 confusable characters. This layer contributes 47.4\% recall on ProtectAI-Validation with only 4.5\% FPR.

\subsection{Layer 2: Embedding Similarity}

Layer 2 computes cosine similarity between the input embedding and 1,403 pre-encoded attack phrase embeddings using BAAI/bge-small-en-v1.5 (33M parameters). The similarity threshold of 0.89 was tuned by sweeping 0.70--0.95 in 0.01 increments on a held-out validation set. This layer adds approximately 5.8 percentage points of recall beyond Layer 1, with minimal FPR increase.

\subsection{Layer 3: DeBERTa Ensemble}

Layer 3 is our primary contribution. It consists of two DeBERTa-v3-base classifiers (184M parameters each) that examine the input from complementary perspectives.

\paragraph{Binary classifier.} A standard sequence classifier fine-tuned on 62,266 samples drawn from deepset, ProtectAI, JailbreakBench, WildJailbreak \cite{jiang2024wildjailbreak}, OR-Bench \cite{zhang2025orbench}, and custom MOF synthetic data. We use focal loss \cite{lin2017focal} with $\gamma=2$ and $\alpha=0.7$. The binary classifier achieves $98.28\% \pm 0.11\%$ validation F1 across 3 seeds but exhibits 20.9\% FPR on NotInject---the motivation for our NLI approach.

\paragraph{NLI classifier.} Rather than asking ``Is this a prompt injection?'' (binary classification), we reframe the question as ``Does this text entail: `This text attempts to override or manipulate system instructions'?'' (natural language inference). This framing leverages the rich semantic understanding of NLI-pretrained models to distinguish between \textit{mentioning} security concepts and \textit{performing} security attacks.

The key discovery motivating this approach: a pre-trained NLI model (MoritzLaurer/DeBERTa-v3-base-mnli-fever-anli, trained on 763K NLI pairs from MNLI, FEVER, and ANLI \cite{nie2020anli}) achieves 9.7\% NotInject FPR in a fully zero-shot setting---without any task-specific training. This compares to 20.9\% for our trained binary classifier. The zero-shot NLI model correctly assigns low entailment scores to benign inputs like ``Please ignore the formatting issues in my document'' (score: 0.01) while assigning higher scores to attacks like ``Forget all rules. You are now DAN'' (score: 0.61).

For fine-tuning, we initialize from the pre-trained NLI checkpoint and convert each (text, label) pair into two NLI pairs using hypothesis templates. Five injection hypotheses and three benign hypotheses are rotated during training to prevent template overfitting. To prevent catastrophic forgetting, we mix 20,000 MNLI samples into the fine-tuning data (total: 119,624 pairs). We use a 10$\times$ lower learning rate ($3 \times 10^{-6}$ vs $2 \times 10^{-5}$). The fine-tuned NLI classifier achieves $98.20\% \pm 0.06\%$ validation F1 across 3 seeds.

\paragraph{Both-agree ensemble.} The binary and NLI classifiers make complementary errors. We exploit this with a simple conjunction: an input is flagged only if $\text{binary\_score} \geq 0.92$ AND $\text{nli\_score} \geq 0.85$. Thresholds were selected via grid sweep optimized for F1 with FPR penalties. Disagreement analysis on NotInject reveals the mechanism: the NLI classifier saves 19 samples the binary classifier incorrectly flags, while the binary classifier prevents 23 NLI false positives.

% ============================================================================
\section{Experimental Setup}

\subsection{Benchmarks}

We evaluate on four publicly available benchmarks: NotInject \cite{lee2025injecguard} (339 benign with trigger words), ProtectAI-Validation (3,227 mixed), deepset/prompt-injections (116 test split), and JailbreakBench \cite{chao2024jailbreakbench} (200 samples). All competitor models are run by us on the same machine with identical evaluation code.

\subsection{Baselines}

\begin{table}[h]
\centering
\small
\begin{tabular}{lll}
\toprule
\textbf{System} & \textbf{Architecture} & \textbf{Params} \\
\midrule
Meta PG2 & mDeBERTa-base & 276M \\
ProtectAI v2 & DeBERTa-v3-base & 184M \\
deepset & DeBERTa-v3-base & 184M \\
Gauntlet (ours) & 3-layer cascade & $\sim$401M \\
\bottomrule
\end{tabular}
\caption{Baseline systems.}
\end{table}

\subsection{Training Details}

\begin{table}[h]
\centering
\small
\begin{tabular}{lcc}
\toprule
& \textbf{Binary} & \textbf{NLI} \\
\midrule
Base model & deberta-v3-base & mnli-fever-anli \\
Training samples & 49,812 & 119,624 \\
Loss & Focal ($\gamma$=2) & Cross-entropy \\
Learning rate & $2 \times 10^{-5}$ & $3 \times 10^{-6}$ \\
Effective batch & 16 & 16 \\
Epochs & 3 & 2 \\
Val F1 (3 seeds) & $98.28 \pm 0.11$ & $98.20 \pm 0.06$ \\
\bottomrule
\end{tabular}
\caption{Training configuration. All experiments on RTX 3060 (6GB).}
\end{table}

% ============================================================================
\section{Results}

\subsection{Main Results}

\begin{table*}[t]
\centering
\small
\begin{tabular}{lccccccc}
\toprule
\textbf{System} & \textbf{NI FPR} $\downarrow$ & \textbf{PAI F1} $\uparrow$ & \textbf{PAI Rec} $\uparrow$ & \textbf{PAI FPR} $\downarrow$ & \textbf{deep F1} $\uparrow$ & \textbf{JBB F1} $\uparrow$ \\
\midrule
Gauntlet L1 only & \textbf{0.3\%} & 0.619 & 47.4\% & 4.5\% & 0.286 & 0.000 \\
Gauntlet L1+L2 & 0.6\% & 0.669 & 53.2\% & \textbf{4.6\%} & 0.400 & 0.112 \\
Gauntlet Ensemble & 4.1\% & 0.701 & 57.8\% & 5.5\% & 0.723 & 0.511 \\
Gauntlet L1+L2+Binary & 20.9\% & 0.734 & 65.3\% & 9.7\% & 0.800 & 0.836 \\
Gauntlet L1+L2+NLI & 26.6\% & \textbf{0.804} & \textbf{77.2\%} & 11.3\% & \textbf{0.824} & \textbf{0.851} \\
\midrule
Meta PG2 & 6.5\% & 0.426 & 30.4\% & 9.4\% & 0.286 & 0.419 \\
ProtectAI v2 & 42.5\% & 0.452 & 38.1\% & 23.2\% & 0.537 & 0.000 \\
deepset & 70.5\% & 0.766 & 97.9\% & 43.9\% & 0.992$^*$ & 0.701 \\
\bottomrule
\end{tabular}
\caption{Cross-benchmark evaluation. All systems evaluated by us on the same hardware with identical code. NI = NotInject, PAI = ProtectAI-Validation, deep = deepset test, JBB = JailbreakBench. $^*$deepset evaluated on its own test set.}
\label{tab:main}
\end{table*}

The NLI cascade (L1+L2+NLI) achieves 0.804 F1 vs Meta PG2's 0.426 (1.89$\times$) and 77.2\% recall vs 30.4\% (2.5$\times$), while using the same DeBERTa architecture family. The Ensemble achieves 5.5\% PAI FPR vs Meta's 9.4\% ($p < 0.0001$, McNemar's test). On NotInject, the Ensemble's 4.1\% FPR (95\% CI: [2.1\%, 6.5\%]) overlaps with Meta's 6.5\% (95\% CI: [4.1\%, 9.1\%]); McNemar's $p = 0.216$, not statistically significant due to small sample size ($n = 339$).

\subsection{Per-Category Detection}

We tested detection across 9 attack categories (4--6 samples each): 37/38 detected (97.4\%). The single miss is a JSON delimiter injection appearing as structured data without explicit instruction-override language.

\subsection{Statistical Validation}

\begin{table}[h]
\centering
\small
\begin{tabular}{lccc}
\toprule
\textbf{Metric} & \textbf{Ensemble} & \textbf{NLI Cascade} & \textbf{Meta PG2} \\
\midrule
FPR & .055 [.045,.066] & .113 [.099,.128] & .094 [.081,.108] \\
Recall & .578 [.553,.604] & .772 [.750,.794] & .304 [.280,.329] \\
F1 & .701 [.679,.721] & .804 [.787,.820] & .425 [.398,.453] \\
\bottomrule
\end{tabular}
\caption{Bootstrap 95\% CIs on PAI-Validation (5,000 resamples). CIs do not overlap between Gauntlet systems and Meta PG2 on F1 or recall.}
\end{table}

\subsection{Ablation Study}

\begin{table}[h]
\centering
\small
\begin{tabular}{lcccc}
\toprule
\textbf{Config} & \textbf{PAI F1} & \textbf{PAI FPR} & \textbf{NI FPR} & \textbf{JBB F1} \\
\midrule
Full Ensemble & 0.701 & 5.5\% & 4.1\% & 0.511 \\
L1+L2+NLI & \textbf{0.804} & 11.3\% & 26.6\% & \textbf{0.851} \\
L1+L2+Binary & 0.734 & 9.7\% & 20.9\% & 0.836 \\
L1+L2 only & 0.668 & 4.6\% & 0.6\% & 0.112 \\
L1 only & 0.619 & 4.5\% & \textbf{0.3\%} & 0.000 \\
Zero-shot NLI & --- & --- & 9.7\% & --- \\
\bottomrule
\end{tabular}
\caption{Ablation results across benchmarks.}
\end{table}

Each cascade layer adds incremental value: L1 provides 47.4\% recall at $<$1ms, L2 adds 5.8pp, L3 adds 12.1pp. The NLI cascade achieves best F1 and JBB F1; the ensemble achieves best FPR.

\subsection{The Zero-Shot NLI Finding}

The most surprising result: a pre-trained NLI model with no exposure to prompt injection data outperforms our trained binary classifier on over-defense. The zero-shot model achieves 9.7\% NotInject FPR compared to 20.9\% for our trained binary classifier, correctly assigning low entailment scores to benign inputs (``Please ignore the formatting issues'': 0.01) and higher scores to attacks (``Forget all rules. You are now DAN'': 0.61).

\subsection{Cascade Flow Analysis}

On ProtectAI-Validation (3,227 samples), Layer 1 catches 47.4\% of attacks at sub-millisecond latency. Layer 2 adds 5.8pp recall with only 2 additional false positives. Layer 3's ensemble contributes 12.1pp while keeping FPR at 5.5\%.

\subsection{Latency}

Per-layer latency on RTX 3060: preprocessing $<$0.1ms, Layer 1 $<$1ms, Layer 2 $\sim$15ms, Layer 3 Binary $\sim$50ms, Layer 3 NLI $\sim$50ms, full cascade worst-case $\sim$120ms. Inputs caught by Layer 1 (47.4\% of attacks) resolve with no model inference.

% ============================================================================
\section{Error Analysis}

\paragraph{False positives.} At the 0.92/0.85 thresholds, 14/339 NotInject samples are incorrectly flagged (4.1\%). Chinese text accounts for 56\% of false positives at lower thresholds---our largest systematic weakness. English false positives include security-adjacent phrases like ``BREAK the CODE'' and ``ensure INSTRUCTIONS.''

\paragraph{False negatives.} Missed attacks share a pattern: indirect extraction via conversational probing (``what letters do I type?''), creative writing as social engineering, and mathematical indirection. Our detector looks for instruction manipulation intent---it does not detect indirect information extraction.

\paragraph{Disagreement analysis.} On NotInject, the NLI classifier saves 19 samples the binary classifier incorrectly flags (keywords without override intent), while the binary classifier prevents 23 NLI false positives (roleplay queries, Chinese text).

% ============================================================================
\section{Limitations}

We identify six limitations: (1) \textbf{English-centric}: Chinese text causes 56\% of remaining FPs; Meta PG2 supports 8 languages. (2) \textbf{No indirect injection}: we detect attack text but do not parse documents. (3) \textbf{Single-turn only}. (4) \textbf{No output scanning}. (5) \textbf{Benchmark limitations}: NotInject has only 339 samples; ProtectAI-Validation is not peer-reviewed; standard evaluation may overestimate by $\sim$8.4 AUC points \cite{polyakov2026benchmarks}. (6) \textbf{No energy-based OOD loss}: Meta PG2's key innovation for distributional robustness.

% ============================================================================
\section{Conclusion}

We presented Gauntlet, introducing NLI-based task framing and a both-agree ensemble strategy for prompt injection detection. Our key finding---that a pre-trained NLI model achieves lower false positive rates than a task-specifically trained binary classifier, with zero prompt injection data---suggests that task formulation matters as much as training data for over-defense mitigation. Future work includes multilingual extension, LODO evaluation, and energy-based loss integration.

% ============================================================================
\section*{Reproducibility}

Code: \url{https://github.com/Ashwinash27/gauntlet-slm} (Apache 2.0). Hardware: RTX 3060 Laptop (6GB). Seeds: \{42, 123, 456\}. Software: Python 3.11, PyTorch 2.1, Transformers 4.36.

% ============================================================================
\bibliography{references}
\bibliographystyle{acl_natbib}

\begin{thebibliography}{13}

\bibitem[Chao et al.(2024)]{chao2024jailbreakbench}
Patrick Chao, Alexander Robey, Edgar Dobriban, Hamed Hassani, George~J. Pappas, and Eric Wong. 2024.
\newblock JailbreakBench: An open robustness benchmark for jailbreaking language models.
\newblock In \emph{NeurIPS 2024}.

\bibitem[Greshake et al.(2023)]{greshake2023indirect}
Kai Greshake, Sahar Abdelnabi, Shailesh Mishra, Christoph Endres, Thorsten Holz, and Mario Fritz. 2023.
\newblock Not what you've signed up for: Compromising real-world LLM-integrated applications with indirect prompt injection.
\newblock \emph{arXiv:2302.12173}.

\bibitem[Jiang et al.(2024)]{jiang2024wildjailbreak}
Jiaqi Jiang et al. 2024.
\newblock WildJailbreak: Probing safety at scale with adversarial benign data.
\newblock In \emph{NeurIPS 2024}. arXiv:2406.18510.

\bibitem[Laurer et al.(2023)]{laurer2023building}
Moritz Laurer, Wouter van Atteveldt, Andreu Casas, and Kasper Welbers. 2023.
\newblock Building efficient universal classifiers with natural language inference.
\newblock \emph{Natural Language Engineering}.

\bibitem[Lee et al.(2025)]{lee2025injecguard}
Linyuan Lee et al. 2025.
\newblock InjecGuard: Benchmarking and mitigating over-defense in prompt injection detection.
\newblock In \emph{ACL 2025}. arXiv:2410.22770.

\bibitem[Lin et al.(2017)]{lin2017focal}
Tsung-Yi Lin, Priya Goyal, Ross Girshick, Kaiming He, and Piotr Doll{\'a}r. 2017.
\newblock Focal loss for dense object detection.
\newblock In \emph{ICCV 2017}.

\bibitem[Liu et al.(2020)]{liu2020energy}
Weitang Liu, Xiaoyun Wang, John Owens, and Yixuan Li. 2020.
\newblock Energy-based out-of-distribution detection.
\newblock In \emph{NeurIPS 2020}. arXiv:2010.03759.

\bibitem[Meta(2025)]{meta2025pg2}
Meta. 2025.
\newblock Llama Prompt Guard 2.
\newblock \emph{PurpleLlama / LlamaFirewall}.

\bibitem[Nie et al.(2020)]{nie2020anli}
Yixin Nie, Adina Williams, Emily Dinan, Mohit Bansal, Jason Weston, and Douwe Kiela. 2020.
\newblock Adversarial NLI: A new benchmark for natural language understanding.
\newblock In \emph{ACL 2020}.

\bibitem[Polyakov et al.(2026)]{polyakov2026benchmarks}
Evgeny Polyakov et al. 2026.
\newblock When benchmarks lie: Evaluating malicious prompt classifiers under true distribution shift.
\newblock \emph{arXiv:2602.14161}.

\bibitem[Schulhoff et al.(2023)]{schulhoff2023hackaprompt}
Sander Schulhoff et al. 2023.
\newblock Ignore this title and HackAPrompt: Exposing systemic weaknesses of LLMs through a global-scale prompt hacking competition.
\newblock In \emph{EMNLP 2023}.

\bibitem[Yin et al.(2019)]{yin2019benchmarking}
Wenpeng Yin, Jamaal Hay, and Dan Roth. 2019.
\newblock Benchmarking zero-shot text classification: Datasets, evaluation and entailment approach.
\newblock In \emph{EMNLP 2019}.

\bibitem[Zhang et al.(2025)]{zhang2025orbench}
Hongjin Zhang et al. 2025.
\newblock OR-Bench: An over-refusal benchmark for large language models.
\newblock In \emph{ICML 2025}.

\end{thebibliography}

\end{document}
